\documentclass[../../main.tex]{subfiles}
\begin{document}
\begin{oframed}
  \subfile{problem}
\end{oframed}

{\bf [解答]} ドモアブルの定理から$a^6=1$だから，$k=0,1,2,\dots,5$に対して
\begin{align}
  a^{6n+k}=a^k\label{eq:1}
\end{align}
となる．

\paragraph{(1)}
\cref{eq:1}から$a^k$についてのみ考えればよいが，このとき$0\le i<j\le5$に対して
$a^i=a^j$と仮定すると$a^{j-i}=1$となって矛盾．故に$a^i\not=a^j$だから$a^k$はすべて
異なり，求めるのは$6$個である．

\paragraph{(2)}
合同式の法を$6$とする．

\subparagraph{(i)$n\equiv \pm1$の時}
\begin{align}
  \{n,2n,3n,4n,5n\}\equiv\{1,2,3,4,5\}
\end{align}
であるから，\cref{eq:1}から(与式)$=1$である．

\subparagraph{(ii)otherwise}

$\{n,2n,3n,4n,5n\}$の中に合同式で$0$に等しいものがあるので，(与式)$=0$となる．
\casesend

以上から
\begin{align}
  \text{(与式)}=\left\{
    \begin{array}{ll}
      1 &  (n\equiv \pm1)     \\
      0 &  (otherwise)
    \end{array}
    \right.\tag{答}
  \end{align}
  となる．
  \newpage
  \end{document}
